% généré par scripts/06_textes.py
\dosmeta{Année de classement : 399 (directe : nom consulaire lu). Références : ICUR I 1453 = ILCV 3432 = CLE 679. Rattachement selon l’EDB : \emph{Incertae originis} (le volume I des ICVR réunit les inscriptions dont la provenance est perdue). Conservation selon l’EDB : \emph{Musei Vaticani: Lapidario Cristiano}. Support : Plaque de marbre. 26 × 32 cm (h. × l.). Inscription gravée. Fiche EDB : Filippo Piazzolla, première version 2015. Notice du catalogue p.~\pageref{not:433}.}
\txtbloc{Texte transmis (EDCS, état 2022)}{\lnum{1}Hic iace⸢t⸣ extinctus pri[mo sub limine vitae]\par
\lnum{2}augusto multum dilec[tus tempore pupus]\par
\lnum{3}nec reor hunc lacrimis [fas esse urgere beatum]\par
\lnum{4}corporis exutus vinc(u)lis q[ui gaudet in astris]\par
\lnum{5}nec mala terreni sensit c[ontagia sensus]\par
\lnum{6}menses namque nove[m postquam super addidit anno]\par
\lnum{7}conplevitque dies vitae s[ex ivit ad astra]\par
\lnum{8}[dep(ositus)] L(ucio) Mallio The(o)doro v(iro) c(larissimo) co(n)s(ule) pri(die) Non[as - -]\par
\lnum{9}q(ui) vixit annum m(enses) VIIII [dies VI]}
\txtbloc{Texte de l’EDB (EDB31243)}{\lnum{1}hic iacet ext̲i̲n̲c̲t̲u̲s̲ p̲r̲i̲[mo sub limine vitae]\par
\lnum{2}angusto mul̲t̲u̲m̲ d̲i̲l̲e̲c̲t̲[us tempore pupus]\par
\lnum{3}nec reor hunc̲ l̲a̲c̲r̲i̲̲m̲̲i̲s̲ [fas esse urgere beatum]\par
\lnum{4}corporis exut̲u̲s̲ v̲i̲n̲c̲l̲i̲s̲ q̲[ui gaudet in astris]\par
\lnum{5}nec mala terr̲e̲n̲i̲ s̲e̲n̲s̲i̲t̲ c̲[ontagia sensus]\par
\lnum{6}menses namq̲u̲e̲ n̲o̲v̲e̲[m postquam super addidit anno]\par
\lnum{7}co⌜m⌝plevitque d̲i̲e̲s̲ v̲i̲t̲a̲e̲ s̲e̲[x ivit ad astra]\par
\lnum{8}[depositus F]l(avio) Mallio Thedoro v(iro) c(larissimo) c̲o̲(n)s̲(ule) p̲r̲i̲(die)  n̲o̲n̲(as) [- - -]\par
\lnum{9}q(ui) vixit an(num) un̲u̲m̲ m̲(enses)  V̲I̲I̲I̲I̲ [d(ies) VI]}
